% ============================================================
%  Template Beamer — Tema Madrid
%  Investigación: Metales traza en muestras biológicas
% ============================================================
\documentclass[aspectratio=169, 11pt]{beamer}

% ── Tema ────────────────────────────────────────────────────
\usetheme{Madrid}
\usefonttheme{professionalfonts}
\useinnertheme{rounded}

% ── Fuentes ─────────────────────────────────────────────────
% Compatible con pdfLaTeX (compilador por defecto en Overleaf)
% Si usas XeLaTeX, reemplaza estas tres líneas por: \usepackage{fontspec}
\usepackage[T1]{fontenc}
% \usepackage[utf8]{inputenc}  % Solo para pdfLaTeX; comentar con XeLaTeX
\usepackage[spanish]{babel}
\usepackage[scaled=0.95]{FiraSans}
\usepackage[scaled=0.85]{FiraMono}
\renewcommand{\familydefault}{\sfdefault}
\linespread{1.15}                 % Más aire entre líneas

% ── Paquetes científicos ─────────────────────────────────────
\usepackage{booktabs}             % Tablas profesionales
\usepackage{siunitx}              % Unidades SI (\SI{}, \si{})
\usepackage{amsmath, amssymb}
\usepackage{graphicx}
% \usepackage{pgfplots}             % Descomentar solo si usas gráficos pgfplots inline
% \pgfplotsset{compat=1.18}


% ── Colores personalizados ───────────────────────────────────
\definecolor{cNegro}     {HTML}{000000}   % Negro
\definecolor{cNavy}      {HTML}{14213D}   % Azul marino oscuro
\definecolor{cDorado}    {HTML}{FCA311}   % Amarillo dorado (acento)
\definecolor{cGrisClaro} {HTML}{E5E5E5}   % Gris claro (fondos)
% \definecolor{cBlanco}  {HTML}{FFFFFF}   % Blanco (usar directamente como 'white')

\setbeamercolor{palette primary}   {bg=cNavy,    fg=white}
\setbeamercolor{palette secondary} {bg=cNegro,   fg=white}
\setbeamercolor{palette tertiary}  {bg=cNavy,    fg=white}
\setbeamercolor{palette quaternary}{bg=cNegro,   fg=white}
\setbeamercolor{structure}         {fg=cNavy}
\setbeamercolor{frametitle}        {bg=cNavy,    fg=white}
\setbeamercolor{title}             {fg=white}
\setbeamercolor{subtitle}          {fg=cDorado}
\setbeamercolor{alerted text}      {fg=cDorado}
\setbeamercolor{example text}      {fg=cNavy}
\setbeamercolor{block title}       {bg=cNavy,    fg=white}
\setbeamercolor{block body}        {bg=cGrisClaro, fg=cNegro}
\setbeamercolor{block title alerted}{bg=cNegro,  fg=cDorado}
\setbeamercolor{block body alerted} {bg=cGrisClaro, fg=cNegro}
\setbeamercolor{block title example}{bg=cDorado, fg=cNegro}
\setbeamercolor{block body example} {bg=cGrisClaro, fg=cNegro}

% ── Ajustes de layout ─────────────────────────────────────────
\setbeamertemplate{navigation symbols}{}  % Quita barra de navegación
\setbeamertemplate{footline}{}            % Quita footline (autor/título/página)
% Descomenta la siguiente línea si quieres solo número de página:
% \setbeamertemplate{footline}[frame number]

% ── Información de la presentación ───────────────────────────
\title{Título de la Presentación}
\subtitle{Subtítulo o nombre del proyecto}
\author{Stefano [Apellido] \and Colaborador Dos}
\institute{%
  Departamento de Química \\
  Universidad / Centro de Investigación \\
  {\small\texttt{correo@institucion.cl}}
}
\date{\today}
% \titlegraphic{\hfill\includegraphics[height=1.2cm]{logo_institucion}}  % Descomentar al agregar logo

% ============================================================
\begin{document}
% ============================================================

% ── Portada ──────────────────────────────────────────────────
\maketitle

% ── Slide de referencia de diseño (eliminar en versión final) ─
\begin{frame}{Referencia de diseño}
  \begin{columns}[T]

    % ── Columna izquierda: paleta ──
    \begin{column}{0.55\textwidth}
      {\large\textbf{Paleta de colores}}\\[0.8em]

      \begin{tabular}{ccccc}
        \colorbox{cNavy}{\rule{0pt}{1.2em}\hspace{2.2em}} &
        \colorbox{cNegro}{\rule{0pt}{1.2em}\hspace{2.2em}} &
        \colorbox{cDorado}{\rule{0pt}{1.2em}\hspace{2.2em}} &
        \colorbox{cGrisClaro}{\rule{0pt}{1.2em}\hspace{2.2em}} &
        \colorbox{white}{\rule{0pt}{1.2em}\hspace{2.2em}} \\[0.3em]
        {\tiny\texttt{\#14213D}} &
        {\tiny\texttt{\#000000}} &
        {\tiny\texttt{\#FCA311}} &
        {\tiny\texttt{\#E5E5E5}} &
        {\tiny\texttt{\#FFFFFF}} \\[0.2em]
        {\tiny Principal} &
        {\tiny Secundario} &
        {\tiny Acento} &
        {\tiny Fondos} &
        {\tiny Blanco} \\
      \end{tabular}
    \end{column}

    % ── Columna derecha: tamaños y estilos ──
    \begin{column}{0.42\textwidth}
      {\large\textbf{Tamaños por elemento}}\\[0.5em]
      {\LARGE Título de slide}\\[0.2em]
      {\Large Subtítulo / bloque}\\[0.2em]
      {\normalsize Texto de cuerpo}\\[0.2em]
      {\small Texto secundario}\\[0.2em]
      {\footnotesize Pie de figura / nota}\\[0.8em]

      {\large\textbf{Estilos}}\\[0.4em]
      \textbf{Negrita}\quad
      \textit{Itálica}\quad
      \texttt{Código}\\[0.3em]
      \textcolor{cDorado}{\textbf{Acento dorado}}\quad
      \alert{\textbf{Alerta}}\\
    \end{column}

  \end{columns}
\end{frame}

% ── Tabla de contenidos ───────────────────────────────────────
\begin{frame}{Contenido}
  \setbeamertemplate{section in toc}[sections numbered]
  \tableofcontents[hideallsubsections]
\end{frame}

% ============================================================
\section{Introducción}
% ============================================================

\begin{frame}{Contexto y motivación}
  \begin{columns}[T]
    \begin{column}{0.55\textwidth}
      \begin{itemize}
        \item Punto 1: contexto general del problema.
        \item Punto 2: relevancia clínica / epidemiológica.
        \item Punto 3: brecha de conocimiento que se aborda.
      \end{itemize}
    \end{column}
    \begin{column}{0.42\textwidth}
      % ↓ Reemplazar con: \includegraphics[width=\linewidth]{tu_figura}
      \fbox{\parbox[c][3.5cm][c]{\linewidth}{\centering\textit{figura\_intro}}}
      \vspace{0.3em}
      {\footnotesize\textit{Pie de figura. Fuente: [Autor, año].}}
    \end{column}
  \end{columns}
\end{frame}

\begin{frame}{Objetivos}
  \begin{block}{Objetivo general}
    Describir en una oración concisa el objetivo principal del estudio.
  \end{block}

  \vspace{0.5em}
  \textbf{Objetivos específicos:}
  \begin{enumerate}
    \item Primer objetivo específico.
    \item Segundo objetivo específico.
    \item Tercer objetivo específico.
  \end{enumerate}
\end{frame}

% ============================================================
\section{Materiales y Métodos}
% ============================================================

\begin{frame}{Diseño del estudio}
  \begin{columns}[T]
    \begin{column}{0.5\textwidth}
      \begin{block}{Población}
        \begin{itemize}
          \item $n = $ \textbf{XX} pacientes
          \item Criterios de inclusión: \ldots
          \item Grupos: Control / Litiasis / Cáncer
        \end{itemize}
      \end{block}
    \end{column}
    \begin{column}{0.5\textwidth}
      \begin{block}{Muestras biológicas}
        \begin{itemize}
          \item Tejido, suero, bilis, sangre
          \item Técnica analítica: ICP-MS / ICP-OES
          \item Metales cuantificados: Cu, Zn, Fe, Se, \ldots
        \end{itemize}
      \end{block}
    \end{column}
  \end{columns}

  \vspace{0.6em}
  \begin{alertblock}{Control de calidad}
    Describir CRM, blancos, duplicados y límites de detección.
  \end{alertblock}
\end{frame}

\begin{frame}{Análisis estadístico}
  \begin{itemize}
    \item Detección/reemplazo de outliers: método IQR de Tukey (mediana grupal).
    \item Prueba de tendencia: \textbf{Jonckheere–Terpstra} (JT test),\\
          orden: Litiasis $<$ Displasia $<$ Cáncer.
    \item Corrección de múltiples comparaciones: \textbf{Benjamini–Hochberg} (FDR).
    \item Umbral de significancia: $\alpha = 0{,}05$ (valor $p$ ajustado).
    \item Software: \textsf{R} v4.x — paquetes \texttt{DescTools}, \texttt{ggplot2}.
  \end{itemize}
\end{frame}

% ============================================================
\section{Resultados}
% ============================================================

\begin{frame}{Concentraciones de metales por grupo}
  % ── Tabla de resultados ──
  \begin{table}
    \centering
    \caption{Mediana (RIC) de concentraciones metálicas (\si{\micro\gram\per\gram}).}
    \small
    \begin{tabular}{lSSS r}
      \toprule
      \textbf{Metal} &
      {\textbf{Litiasis}} &
      {\textbf{Displasia}} &
      {\textbf{Cáncer}} &
      \textbf{$p_\text{adj}$} \\
      \midrule
      Cu   & 1.23 & 2.10 & 4.56 & $<0.001$ \\
      Zn   & 8.90 & 7.50 & 5.30 & 0.012    \\
      Fe   & 45.2 & 50.1 & 62.4 & 0.034    \\
      Se   & 0.45 & 0.40 & 0.31 & 0.087    \\
      \bottomrule
    \end{tabular}
  \end{table}
  {\footnotesize JT test, 5\,000 permutaciones; corrección BH.}
\end{frame}

\begin{frame}{Distribución de Cu y Zn — Boxplots}
  \begin{columns}[T]
    \begin{column}{0.49\textwidth}
      % ↓ Reemplazar con: \includegraphics[width=\linewidth]{boxplot_Cu}
      \fbox{\parbox[c][4cm][c]{\linewidth}{\centering\textit{boxplot\_Cu}}}
      \begin{center}
        {\small Cobre (Cu)}
      \end{center}
    \end{column}
    \begin{column}{0.49\textwidth}
      % ↓ Reemplazar con: \includegraphics[width=\linewidth]{boxplot_Zn}
      \fbox{\parbox[c][4cm][c]{\linewidth}{\centering\textit{boxplot\_Zn}}}
      \begin{center}
        {\small Zinc (Zn)}
      \end{center}
    \end{column}
  \end{columns}
  \vspace{0.3em}
  {\footnotesize $^\ast p_\text{adj} < 0{,}05$;\quad $^{\ast\ast} p_\text{adj} < 0{,}01$;\quad ns = no significativo.}
\end{frame}

\begin{frame}{Razones metálicas}
  \begin{columns}[T]
    \begin{column}{0.48\textwidth}
      % ↓ Reemplazar con: \includegraphics[width=\linewidth]{boxplot_CuZn_ratio}
      \fbox{\parbox[c][4cm][c]{\linewidth}{\centering\textit{boxplot\_CuZn\_ratio}}}
      \begin{center}
        {\small Razón Cu/Zn}
      \end{center}
    \end{column}
    \begin{column}{0.48\textwidth}
      \begin{exampleblock}{Interpretación}
        \begin{itemize}
          \item La razón Cu/Zn aumenta progresivamente de litiasis a cáncer.
          \item Consistente con desequilibrio redox en tejido tumoral.
          \item \textit{Agregar más contexto según tus datos.}
        \end{itemize}
      \end{exampleblock}
    \end{column}
  \end{columns}
\end{frame}

% ============================================================
\section{Discusión y Conclusiones}
% ============================================================

\begin{frame}{Discusión}
  \begin{itemize}
    \item \textbf{Hallazgo principal:} \ldots\ concuerda/difiere con [Autor, año].
    \item Posible mecanismo biológico que explica el patrón observado.
    \item Limitaciones: tamaño muestral, diseño transversal, \ldots
    \item Fortalezas: control de calidad analítico, diseño histológico confirmado.
  \end{itemize}
\end{frame}

\begin{frame}{Conclusiones}
  \begin{enumerate}
    \item Primera conclusión derivada directamente de los resultados.
    \item Segunda conclusión.
    \item Implicación clínica o proyección hacia trabajo futuro.
  \end{enumerate}

  \vspace{1em}
  \begin{alertblock}{Próximos pasos}
    Indicar el siguiente experimento, análisis o publicación planificada.
  \end{alertblock}
\end{frame}

% ── Agradecimientos ───────────────────────────────────────────
\begin{frame}{Agradecimientos}
  \begin{itemize}
    \item Nombre de financiamiento / proyecto / beca.
    \item Colaboradores: Ignacio, Lizeth, Gina y equipo clínico.
    \item Institución / hospital participante.
  \end{itemize}
\end{frame}


% ── Apéndice (slides extra, no contabilizadas) ────────────────
\appendix

\begin{frame}{Apéndice — Datos adicionales}
  Diapositivas de respaldo para preguntas del público.\\
  No se muestran en la presentación principal.
\end{frame}

\end{document}
